\section{A full gap for arbitrary two-level potentials}\label{twolevelgap}
The stationary floor and critical cap yield a stronger conclusion when the
potential has at most two values. No marginal matching is needed in this
section. The level ratio and the shape or mass of the high-level set may
vary with dimension. The restriction to two values is explicit; no reduction
of general potentials to this class is asserted.

\begin{theorem}\label{twolevelfullgap}
Fix $0<\rho<1$ and $0<b<1$. Let $h>0$ take at most two values on
$\{-1,1\}^L$, $L\geq1$, and normalize $\E hP_\rho h=1$. Suppose
$\pi=hP_\rho h\geq b$ and $\E\pi^s\leq2$, where
$s=(1+\rho)/(1-\rho)$. The actual Doob kernel $K$ has
\[
 \operatorname{gap}(K)\geq G(\rho,b)>0
\]
uniformly in $L$ and $h$. One admissible, deliberately conservative choice is
\begin{align*}
 s'&=s/(s-1), & B&=2^{1/s}, & \kappa&=4B/b^2,\\
 \eta&=\frac{1-2^{-1/s}}{1+\kappa}, &
 t_*&=\min\left\{2^{-s},\left(\frac{\eta^2}{4\kappa}\right)^s\right\}, &
 \epsilon_*&=bt_*,\\
 \gamma_*&=\frac{b^2(1-\rho)}{4\,2^{s'}B^{2s'}}
                 \epsilon_*^{s'-1}, &
 \Phi_*&=\min\{\eta/2,\gamma_*\}, &
 G(\rho,b)&=\min\left\{\frac{b^2(1-\rho)}{16},\frac{\Phi_*^2}{4}\right\}.
\end{align*}
All symbols in these constants are local to this section. No optimality is
claimed.
\end{theorem}
\begin{proof}
A constant potential normalizes to $h=1$ and gives $K=P_\rho$, whose gap is
$1-\rho$. Otherwise write $h=M$ on $A$ of mass $\alpha\in(0,1)$ and $h=m$
on $A^c$, where $0<m<M$. Put $P=P_\rho$, $g=Ph$, and $\nu=\pi\mu$.

\paragraph{A lower bound derived from the floor.}
Set $a=\E h=\E g$. The constant Fourier component gives
$1=\|P_{\sqrt\rho}h\|_2^2\geq a^2$, so $a\leq1$.
Integrating $mg\geq b$ on $A^c$ gives
$m\E(g1_{A^c})\geq b(1-\alpha)$, and hence $m\geq b(1-\alpha)$.
Since $g\leq M$, the floor on $A^c$ also gives $mM\geq b$.
The bound $M\alpha\leq a\leq1$ then yields $m\geq b/M\geq b\alpha$.
Consequently
\begin{equation}\label{twolevelminimum}
 m\geq b\max\{\alpha,1-\alpha\}\geq b/2.
\end{equation}
In particular $g\geq m$, so on $A$ one has $\pi=Mg\geq Mm$.
The critical cap implies $(Mm)^s\alpha\leq2$, or
\begin{equation}\label{twolevelbalance}
 \frac M m\alpha^{1/s}\leq\frac B{m^2}\leq\kappa.
\end{equation}
This controls a mass-weighted ratio; $M/m$ itself remains unrestricted.

\paragraph{Reference restrictions and the high-level denominator.}
For $Q=P_{\sqrt\rho}$, ordinary hypercontractivity gives
$\|Qf\|_2\leq\|f\|_{1+\rho}$~\cite{odonnelllecture}.
If $f$ is supported on $U$, H\"older's inequality gives
$\|f\|_{1+\rho}\leq\mu(U)^{1/(2s)}\|f\|_2$.
Using $P=Q^2$ in the bilinear form therefore proves, for all sets $U,V$,
\begin{equation}\label{twolevelrestriction}
 \|1_U P1_V\|_{2\to2}\leq[\mu(U)\mu(V)]^{1/(2s)}.
\end{equation}
The norms in this paragraph use $L^2(\mu)$.

Let $u=h/g$ and $S=M_{\sqrt u}PM_{\sqrt u}$. Multiplication by $\sqrt\pi$
is an isometry from $L^2(\nu)$ to $L^2(\mu)$ conjugating $K$ to $S$.
It commutes with restrictions, and $S$ is positive definite.
On $A$, define $P_A=1_A P1_A$, $R=\operatorname{diag}(P1_A)|_A$,
$r_0=m/M$, and $q=\alpha^{1/s}$. The internal Dirichlet form gives
$P_A\preceq R$, while \eqref{twolevelrestriction} gives $P_A\preceq qI$.
The exact two-level identity
\[
 g/M=r_0+(1-r_0)P1_A
\]
implies
\[
 S_{AA}=D^{-1/2}P_A D^{-1/2},\qquad D=r_0I+(1-r_0)R.
\]
For $H_0=r_0+(1-r_0)q$, put $\theta=r_0/H_0$ and $\ell=q/H_0$.
Then $0<\theta<1$, $\theta q=\ell r_0$, and
$1-\theta=\ell(1-r_0)$. A convex combination of the two Loewner bounds gives
\[
 P_A\preceq\theta qI+(1-\theta)R=\ell D.
\]
No commutativity between $P_A$ and $R$ is required. Congruence by
$D^{-1/2}$ proves
\begin{equation}\label{twolevelkilled}
 1-\|S_{AA}\|\geq1-\ell
 =\frac{1-q}{1+(M/m-1)q}.
\end{equation}
Assume first that $\alpha\leq1/2$. By \eqref{twolevelbalance}, the last
expression is at least $\eta$. Every compression $S_{UU}$ with $U\subset A$
then has norm at most $1-\eta$.

\paragraph{Small cuts, including cuts meeting both levels.}
For $0<\nu(C)\leq\epsilon_*$, put $U=C\cap A$ and $V=C\cap A^c$.
The floor gives $\mu(V)\leq\nu(C)/b\leq t_*$.
On $A^c$ the multiplier satisfies $u=m/g\leq1$; on $A$ it satisfies
$u=M/g\leq M/m$. Thus \eqref{twolevelrestriction} gives
\[
 \|S_{VV}\|\leq\mu(V)^{1/s}\leq\frac12,
\]
and, using \eqref{twolevelbalance},
\[
 \|S_{UV}\|\leq\sqrt{M/m}\,[\mu(U)\mu(V)]^{1/(2s)}
 \leq\sqrt\kappa\,\mu(V)^{1/(2s)}\leq\eta/2.
\]
Since $0<\eta<1/2$, the block-diagonal part on $U\oplus V$ has norm at
most $1-\eta$. The self-adjoint off-diagonal block has norm
$\|S_{UV}\|$. The triangle inequality gives
$\|S_{CC}\|\leq1-\eta/2$, also when one part is empty.

Define the actual stationary outward flow
\[
 F_h(C)=\sum_{x\in C,\,y\notin C}\mu(x)P(x,y)h(x)h(y)
       =\mathcal D_K(1_C).
\]
The vector $\sqrt\pi\,1_C$ has squared $L^2(\mu)$ norm $\nu(C)$.
Conjugacy and the bound on $S_{CC}$ yield
\begin{equation}\label{twolevelsmallflow}
 F_h(C)=\nu(C)-\langle1_C,K1_C\rangle_\nu
       \geq(\eta/2)\nu(C).
\end{equation}

\paragraph{The remaining cuts and Cheeger's inequality.}
Since $h\geq m$, the actual edge measure gives
$\mathcal D_K\geq m^2\mathcal D_P$. Hence
\[
 F_h(C)\geq m^2(1-\rho)\mu(C)\mu(C^c).
\]
The cap supplies $\nu(E)\leq B\mu(E)^{1/s'}$ for every $E$.
For $\epsilon_*\leq\nu(C)\leq1/2$, it follows that
\[
 \mu(C)\geq(\nu(C)/B)^{s'},\qquad
 \mu(C^c)\geq(1/(2B))^{s'}.
\]
Combining these estimates with \eqref{twolevelminimum} gives
\begin{equation}\label{twolevelbalancedflow}
 \frac{F_h(C)}{\nu(C)}
 \geq\frac{b^2(1-\rho)}{4\,2^{s'}B^{2s'}}\nu(C)^{s'-1}
 \geq\gamma_*.
\end{equation}
Equations~\eqref{twolevelsmallflow} and \eqref{twolevelbalancedflow}
cover every set relevant to the conductance
\[
 \Phi(K)=\min_{0<\nu(C)\leq1/2}\frac{F_h(C)}{\nu(C)},
 \qquad \Phi(K)\geq\Phi_*.
\]
Use the lazy reversible Cheeger bound from~\cite{peres} on
$\overline K=(I+K)/2$. Its conductance is $\Phi(K)/2$ and its gap is
$\operatorname{gap}(K)/2$. Thus
\[
 \frac{\operatorname{gap}(K)}2
 \geq\frac12\left(\frac{\Phi(K)}2\right)^2,
 \qquad \operatorname{gap}(K)\geq\Phi_*^2/4.
\]
This is deliberately conservative and needs only the cited lazy version.

\paragraph{A high-level set of mass above one half.}
If instead $\alpha>1/2$, then $M\alpha\leq1$ gives $M<2$, whereas
\eqref{twolevelminimum} gives $m\geq b/2$. Since $\pi=hg\leq M^2$,
centering at $\E_\mu f$ shows
$\operatorname{Var}_\nu(f)\leq M^2\operatorname{Var}_\mu(f)$.
The edge comparison gives
$\mathcal D_K(f)\geq m^2(1-\rho)\operatorname{Var}_\mu(f)$.
Taking Rayleigh quotients proves
\[
 \operatorname{gap}(K)\geq(m/M)^2(1-\rho)
                      \geq b^2(1-\rho)/16.
\]
Together with the constant case and the $\alpha\leq1/2$ case this proves
the stated $G(\rho,b)$.
\end{proof}

Theorem~\ref{twolevelfullgap} controls the full chain, not just a quotient
or a killed restriction. It excludes collapsing gaps from arbitrary
constant-height plateaux above a constant positive background under the
original stationary floor and cap. Matching is unnecessary within this
class. The proof derives \eqref{twolevelminimum} from its two-level
structure; it does not establish that bound for arbitrary positive
potentials or replace the original open question by an extra assumption.
